\documentclass[11pt]{article}

\usepackage[margin=1in]{geometry}
\usepackage[T1]{fontenc}
\usepackage[utf8]{inputenc}
\usepackage{lmodern}
\usepackage{setspace}
\usepackage{titlesec}
\usepackage{hyperref}

\setstretch{1.1}

\titleformat{\section}{\large\bfseries}{\thesection}{1em}{}

\begin{document}

\begin{center}
    {\LARGE Final Project Proposal}\\[1.5em]
\end{center}

\begin{tabular}{@{}ll@{}}
    \textbf{Class:} & COMP 5300 \\[0.5em]
    \textbf{Name:} & Jegannarayanan Sivakasiulaganathan \\[0.5em]
    \textbf{Student ID:} & 02197006 \\[0.5em]
    \textbf{Email:} & Jegannarayanan\_Sivakasiulaganathan@student.uml.edu
\end{tabular}

\vspace{2em}

\section{Overview}
My final project proposal is to create a Python AST to LLVM MLIR IR Object
pipeline. The idea is to leverage the library mlir-hs, which consists of
API calls to create MLIR IR objects.
Reference:
https://github.com/google/mlir-hs

\section{Motivation}
The motivation for this project is as follows:

Using Haskell to reason about Python AST objects allows us to leverage
Haskell's early type safety, complex logic expressivity, and reliability to
directly compile Python code.
We can directly work with Python AST nodes, rather than having to go through
CPython/other LLVM workflows. This allows direct lowering from high level syntax (Python)
to MLIR. Note we can operate directly on
function annotations and library calls.

\section{Explanation}
A more concrete explanation of this approach is:
This project will utilize an existing Python AST library (One which takes in Python code and outputs a Python AST).
Given an AST, I'd like to walk it, then call the appropriate mlir-hs functions to
create MLIR-IR objects. Finally, I'll print the ensuing IR. The code I'm writing will call mlir-hs itself.

\section{Scope/Conclusion}
I've scoped this project to act on basic vector ops - multiplication, addition, etc.
The code will essentially be a mapping between those ops and the appropriate MLIR IR objects.
Furthermore, the project will not focus on execution correctness or optimization.
Some ways to expand this scope are to act on tensors, actually have it run and test optimization,
etc.



\end{document}
